\chapter{Literature Review}
\label{chap:literature}
This section reviews the relevant literature about applications of deep probabilistic models in three aspects: Bayesian inference, especially parameter identification; system dynamics modelling and state estimation; data-driven control. The deep probabilistic models mainly include \acp{DGP}, \acp{BNN}, \acp{CVAE}, \acp{CNF}, and also conditional diffusion models. 

\section{Parameter Identification}  
\subsection{Deep Gaussian Processes}
\acfp{DGP} are sequential compositions of Gaussian processes, where the outputs of one Gaussian process are the inputs to the next. They extend the single-layer \ac{GPR} and can represent compositional or hierarchical structure of variables. A survey \citep{jakkala2021survey} summarizes the improved methods of \acp{DGP}.

\citeauthor{abraham2023deepgp} \citeyearpar{abraham2023deepgp} proposed a \ac{DGP}-based method for Bayesian inference of compositional parameters from \acp{PDE}. The parameters are assigned with a \ac{DGP} prior which models compositional functions. Then, the posterior distribution of the parameters is optimized via variational inference.

Applications of \acp{DGP} for Bayesian inference span multiple domains. For instance, \citeauthor{dgp2018park} \citeyearpar{dgp2018park} utilized \acp{DGP} to identify contaminant release sources, while \citeauthor{dgp2023desai} \citeyearpar{dgp2023desai} employed them for air quality estimation. 

In biomedical contexts, \citeauthor{NIPS2017_861dc9bd} \citeyearpar{NIPS2017_861dc9bd} leveraged \acp{DGP} to assess patient survival under competing risks. 

These studies show promising results in comparison to \acp{GPR} as well as other machine learning methods, and highlight the versatility and effectiveness of \acp{DGP} in addressing parameter inference challenges. However, the computational cost and storage cost remain the main drawbacks of \acp{DGP}, especially for high-dimensional data.

\subsection{Bayesian Neural Networks}
Applying \acfp{BNN} for Bayesian inference is usually combined with other methods for better extracting information from observation data. For example, \citeauthor{YANG2021109913} \citeyearpar{YANG2021109913} integrated a \ac{BNN} with \acfp{PINN} (B-PINN) to infer parameters in \acfp{PDE} with Gaussian noise. \acp{PINN} construct a neural network to approximate the forward solution to the underlying \ac{ODE} or \ac{PDE} problem by minimizing reconstruction loss of data, while regularized by the underlying governing \ac{ODE} or \ac{PDE} residual loss. Meanwhile, model parameters can also be estimated probabilistically via \ac{MCMC} or variational inference. The B-PINN achieved more accurate results with stronger noise in observations. 

\citeauthor{Merkatas2023_sm} \citeyearpar{Merkatas2023_sm} used an autoregressive \ac{BNN} for estimating the system parameters, together with a nonparametric noise model to capture complex noise characteristics in the data. The \ac{BNN} network parameters are approximated through \ac{HMC} sampling approach.  \acp{BNN} are flexible to be integrated with other models and can provide uncertainty quantification, at the cost of higher computational complexity. On the other hand, applying \acp{BNN} for Bayesian inference is not straightforward and must be combined with other architectures. Besides, the uncertainty of network weights does not necessarily translate to the uncertainty of the target variables.

\subsection{Conditional Variational Autoencoders}
\acfp{CVAE} represent a conditional extension of the \acf{VAE} framework. \acp{CVAE} have been widely adopted for amortized Bayesian inference across various scientific domains. For example, \citeauthor{vae_sbi} \citeyearpar{vae_sbi} introduced a \ac{CVAE}-based approach for efficient Bayesian inference in scientific simulators.

In astronomy, \acp{CVAE} have enabled Bayesian parameter estimation from gravitational wave data, as demonstrated by \citeauthor{Yamamoto_2020rse} \citeyearpar{Yamamoto_2020rse}, \citeauthor{Gabbard2022} \citeyearpar{Gabbard2022}, and \citeauthor{Sun_2025} \citeyearpar{Sun_2025}. Additionally, \citeauthor{Ferreira_2025} \citeyearpar{Ferreira_2025} applied \acp{CVAE} to infer the equation of state of neutron stars. 

In biomedical engineering, \acp{CVAE} have been used for neuroimaging analysis under varying correlated parameters \citeyearpar{cvae2025}, and for parameter estimation in mouse cardiac muscle studies \citeyearpar{TUNE2025179}. 

These examples highlight the versatility and effectiveness of \acp{CVAE} for Bayesian inference of high-dimensional data. One shortcoming exhibited in the research \citep{Sun_2025} is that the \ac{CVAE} tends to estimate broader posterior distributions compared to other methods, indicating a potential limitation in capturing precise uncertainty.

\subsection{Conditional Normalizing Flows}
BayesFlow (in Section \ref{sec:condnf}) \citep{radev2020bayesflow} provides a general framework for amortized Bayesian inference using \acp{CNF}. Besides our work, it has also been applied to various scientific domains. For instance, \citeauthor{Kucharsk2025} \citeyearpar{Kucharsk2025} applied BayesFlow in multiple cognitive models. \citeauthor{Zeng2023} \citeyearpar{Zeng2023} demonstrated its effectiveness in damage detection in structural health monitoring. \citeauthor{Rahmandad2025_uh} \citeyearpar{Rahmandad2025_uh} utilized BayesFlow to infer parameters from system dynamics, such as the random walk model. 

Some other works share similar idea with BayesFlow. \citeauthor{De_Santi_2024} \citeyearpar{De_Santi_2024} referred the summary network as an embedding neural network. \citeauthor{zhdanov2022gisaxs} \citeyearpar{zhdanov2022gisaxs} extended the condition with richer synthetic data, which is generated by a \ac{CVAE}. \citeauthor{DONG2025117457} \citeyearpar{DONG2025117457} arranged a \ac{CNF} through taking the conditioning variable as an additional input to each coupling layer and enlarging the latent space as well. 

These works demonstrate the effectiveness of \acp{CNF} for amortized Bayesian inference. Especially, the design of the summary network needs to be tailored to the specific characteristics of the data, which is crucial for extracting informative representations from observations and inferring parameters precisely.

\subsection{Conditional Diffusion Models}
Diffusion models have recently emerged as powerful tools for Bayesian inference due to their ability to model complex, multimodal distributions. \citeauthor{zhu2025diffusion} \citeyearpar{zhu2025diffusion} introduced a diffusion-based framework for parameter identification in dynamic power systems, demonstrating improved accuracy over traditional methods. 

Conditional diffusion models have also been applied to infer parameters, unobservable states, and system dynamics in \acp{PDE} \citep{haitsiukevich2024}. 

Furthermore, \citeauthor{baldassari2023} \citeyearpar{baldassari2023} developed a framework for conditional score-based diffusion models in infinite-dimensional function spaces to perform Bayesian inference for linear inverse problems. 

\citeauthor{chen2025conddiff} \citeyearpar{chen2025conddiff} proposed a conditional diffusion model for amortized posterior estimation, showing superior performance compared to \acp{CNF}, particularly in reconstructing multimodal distributions and capturing variable boundaries. 
Nevertheless, the computational cost of inference with diffusion models is much higher than other \acp{DGM} due to its iterative sampling procedure.

\section{System Dynamics Modelling and State Estimation}
After introducing versatile applications of deep probabilistic models for Bayesian inference, mainly parameter identification, this section focuses on the applications of deep probabilistic approaches specific for system dynamics modelling and state estimation. Probabilistic dynamics modelling and state estimation overlap with \acp{SSM} and stochastic processes modelling approaches. 
\subsection{System Dynamics Modelling}
Identification of nonlinear system dynamics solely from input-output data is a challenging task with conventional methods. In contrast, deep probabilistic models, especially \acp{DGM}, have shown great potential in capturing complex nonlinear relationships while providing uncertainty quantification. 

\citeauthor{8972113} \citeyearpar{8972113} stacked \ac{LSTM} layers with a probabilistic layer, obtaining a Bayesian inference model for nonlinear system identification, capturing nonlinear dynamics while providing uncertainty quantification. 

\citeauthor{PANIAGUA2024200344} \citeyearpar{PANIAGUA2024200344} introduced a framework integrating a \ac{VAE} with \ac{NARX} for identifying nonlinear systems using input/output data. The \ac{VAE} encodes the input/output data into a latent space, and the \ac{NARX} decodes the latent space information and reconstructs the system outputs. This approach effectively captures nonlinear relationships effectively and improves the robustness through the \ac{VAE} in comparison to the \ac{NARX} alone. 

\citeauthor{roeder2019abi} \citeyearpar{roeder2019abi} proposed an amortized Bayesian inference method for learning hierarchical nonlinear dynamical systems. The method extends the \ac{VAE} by incorporating the structural information of the system as hierarchical conditions. The encoder performs amortized inference over latent variables conditioned on observed trajectories and system structure. The decoder employs differentiable \acp{ODE} to describe the system dynamics, which can include both known mechanistic components and learned neural components. This method enables integrating knowledge of system structure with data-driven learning, improving interpretability and generalization.

\subsection{Data-Assisted Bayesian Filters}
From the section \ref{sec:bayesfilter} we have known that estimating state from the noisy observations is strongly dependent on the completeness of system models and strong assumptions on noise. Under model mismatches or incompleteness, deep probabilistic alternatives have been proposed to integrate Bayesian filters with offline data. The data-assisted Bayesian filters take advantage of the rigor of the state prediction and update, achieving explainability and better results than solely data-driven methods.

For instance, \citeauthor{yan2024brnn} \citeyearpar{yan2024brnn} introduced a Bayesian gated recurrent framework for non‑Markov dynamical systems, which consists of a memory update gate for capturing the temporal regularities in state evolution, a state prediction gate with the evolution mismatch compensation, and a state update gate with the observation mismatch compensation; it achieved stronger state estimation performance than standard \acp{RNN} and Transformer baselines. 

\citeauthor{fraccaro2017kvae} \citeyearpar{fraccaro2017kvae} proposed the Kalman \ac{VAE}, which combines a \ac{VAE} with a linear Gaussian state-space model. The \ac{VAE} encoder learns a low-dimensional latent representation of the observations, while the Kalman filter performs recursive Bayesian state estimation in the latent space. The decoder reconstructs the observations from filtered latent states, allowing uncertainty-aware tracking of temporal dynamics. Similar ideas have been explored in \citep{pasmans2025ekfl, ZHANG2023106424}. 

\ac{NF} have also been explored to be combined with linear Gaussian state-space model. \citeauthor{wang2025flow} \citeyearpar{wang2025flow} utilized a \ac{NF} to transform nonlinear model with non-Gaussian noise into latent space, where a Kalman filter is operated. The major advantage in \acp{NF} compared with \acp{VAE} lies in the invertibility, which allows efficient sampling of target filtering distribution in the original space. The authors shown that the flow-based Bayesian filter achieves better accuracy and computational efficiency compared to conventional filtering methods in challenging high-dimensional settings.

Another way of applying \ac{NF} is proposed by \citeauthor{chen2021dpfcnf} \citeyearpar{chen2021dpfcnf}. A \ac{CNF} is employed to construct flexible proposal distributions in differentiable particle filters. The \ac{NF} is conditioned on the observed data to migrate particles to areas closer to the true posterior. The integration of \acp{CNF} improved the accuracy in a visual tracking task. 

\subsection{Deep State-Space Models}
Constructing the state transition and observation models in \acp{SSM} using neural networks allows for the representation of a wide range of dynamic behaviors, thanks to their flexibility and expressive capability. The review \citep{lin2024ssm} provides a selective review of recent advancements in deep neural network-based approaches for \acp{SSM}, and presents a unified perspective for discrete time deep state space models and continuous time ones such as latent neural \acp{ODE} and neural \acp{SDE}. Recent advancements in sequence modeling have led to the emergence of structured \acp{SSM} as an efficient alternative to \acp{RNN} and Transformers, addressing challenges in long-range dependency modeling and computational efficiency. A comprehensive survey about structured \acp{SSM} is conducted by \citeauthor{somvanshi2025s4mamba} \citeyearpar{somvanshi2025s4mamba}. 

\acp{NF} are also investigated to model noises in the \acp{SSM}. \citeauthor{lin2025} \citeyearpar{lin2025} proposed an efficient transformed Gaussian process \ac{SSM} that integrates a single shared GP with input-dependent \acp{NF}, yielding an expressive implicit process prior that captures complex, non-stationary transition dynamics while significantly reducing model complexity.

\subsection{Normalizing Flows and Stochastic Processes}
Speaking of probabilistic modelling of time series, different methods have been investigated. In this section, we focus on the research about leveraging \acp{NF} to model stochastic processes. 

\ac{SCNF} \citep{hodg2021stoch} treats \acp{SDE} as random \acp{ODE}. Random \acp{ODE} are a type of differential equation where the vector field includes a stochastic process, which is approximated by deterministic functions. Random \acp{ODE} are trained in the manner of continuous normalizing flows. However, \ac{SCNF} suffers from the overfitting issue of the vanishing diffusion term. 

\ac{CTFP} \citep{deng2021ctfp} models stochastic processes by transforming latent Wiener processes via continuous normalizing flows. The continuous normalizing flows are trained in the forward direction through mapping the data samples to the latent variables. Time-Changed Normalizing Flow \citep{Bekri_2024} extends the base Wiener process in \ac{CTFP} with a time function. Continuous Latent Process Flow \citep{deng2021continuous}, as an improvement of \ac{CTFP}, applies a LatentSDE \citep{li2020scalable} as the encoder to generate latent variables, and a continuous normalizing flow as the decoder. Although these methods shown good performance in generating probabilistic time series, they are computationally expensive and not able to solve initial value problems directly or model dynamical systems with control inputs, which makes them unsuitable for many engineering applications.

Although several works have demonstrated the potential of normalizing flows for modelling stochastic processes, their application to system dynamics modelling and state estimation remains limited. Key challenges include high computational cost, difficulty handling control inputs and initial‑value problems, and risks of overfitting or vanishing diffusion. Addressing these issues represents a promising direction for our research.

\section{Deep Probabilistic Models and Control}
The identified system model can be further utilized in \ac{MPC}, where the modelled uncertainty plays a crucial role in enabling robust decision-making. Moreover, deep probabilistic approaches for capturing input-output relationships are well-suited for integration into model-free \ac{RL} frameworks. Alternatively, \acp{SSM} can be implemented within model-based \ac{RL}, providing structured representations that facilitate planning and control under uncertainty. A survey \citep{zhu2024diffusion} provides a comprehensive overview of recent research efforts focused on the application of diffusion models within the domain of reinforcement learning. Especially, \citeauthor{ghugare2025nf} \citeyearpar{ghugare2025nf} brought up that the probabilistic modelling and inference capability of \acp{NF} suits better for \ac{RL} framework in comparison to diffusion models, including acting as a policy, as a Q-function or value model, and for modeling state densities. 

We summarize \ac{RL} methods incorporating \acp{BNN}, \acp{VAE} and \acp{NF} in Table.